%!TEX root = ../Thesis.tex


\chapter{Project Description}

\section{Background}

When the British archaeologist Arthur Evans in the beginning of the 20th century went to excavate Knossos on Crete, he found a huge amount of clay tablets written in a previously unknown linear script.
Some documents were clearly of a newer script, which he called \textit{Linear B}, while others were of an older one, which he named \textit{Linear A}.
Some fifty years later, the British amateur archaeologist Michael Ventris finally succeeded in deciphering the newer script Linear B, which he proved to be representing Mycenaean Greek, the earliest attested form of Greek.
But the same hypothesis did not hold for the older script Linear A, whose encoded language remained unknown, still yet to be deciphered.

A key obstacle to its decipherment has been the small size of its corpus, lately consisting of only around 2,500 inscriptions (about 1,400 if only including those well-preserved enough to be readable) - significantly smaller than that of its younger counterpart that now spans more than 4,500 inscriptions (which are comparatively longer and better preserved).\footcite{Salgarella2020}\footcite[p.~13,46]{Salgarella2025} 
But new inscriptions are still being discovered, with the corpus most recently extended in 2025 by over 100 new documents.\footcite{RILA-S1}\footcite{CNRLinearACorpus2025}

Another challenge has been the lack of digital representations of these inscriptions, in a form suitable for statistical analysis.
While photographs with accompanying drawings of nearly the entire corpus are available online through \acf{Cefael} (see \textcite{EtudesCretoises21}), these are only available in print form.
This motivated the development of the palæographical\footnote{The term "palæographical" refers to the study of ancient writing systems.} database \textit{The Signs of Linear A} (SigLA), which seeks to "[develop] a systematic, exhaustive and user-friendly open access database of all Linear A inscriptions [...] in order to carry out statistical and palæographic analyses within the epigraphic corpus" \parencite{SalgarellaCastellanSigLA2020}.
However, the last addition to SigLA was in 2021, and the project has so far managed to document just 772 of all currently known inscriptions.\footcite{SigLAAboutPage}\footcite{SigLASearchDocuments}

\subsection{A problem of manual digitization}

The drawings currently stored on SigLA were made by Ester Salgarella using the painting tool Krita.
For each document, she manually created her own annotated drawings based on the originals from \acs{Cefael}.
From these, she proceeded to create layered Krita files with metadata for each sign, which could then at last be imported to SigLA.\footcite{SalgarellaCastellanSigLA2020}
According to her (personal communication), this has been a highly time-consuming undertaking. 
And the necessity of this kind of labor-intensive work has arguably been the biggest reason to why SigLA still does not contain the entire available corpus.

The aim of this project is therefore to investigate whether parts of this process can be automated. More specifically, whether some kind of semi-automatic segmentation tool 
can accelerate the workflow, while maintaining high-quality outputs that are still compatible with SigLA's database.

\subsection{A problem of small data}
The small size of the Linear A corpus and the irregularities in quality of its documents is a huge problem for the application of deep learning approaches, which typically require large amounts of data to perform well.
Consequently, this project will instead seek to make use of generalized machine learning and computer vision, in particular Meta's state-of-the-art interactive segmentation tool, the \acf{SAM} \parencite{sam}, which is pre-trained on a massive and diverse dataset of over 1 billion masks.
\acs{SAM}'s generalization and interactivity (allowing user-placed markers that guide the model) might thus allow for semi-automatic annotation of Linear A inscriptions without the need for large amounts of training data, which is unfortunately not available in this context.\footnote{By the "interactivity" of \acs{SAM}, what is meant is NOT any built-in user interface in the form of a GUI, but rather the possibility of prompt-based placement of rectangles of focus, and points of additions and deletions that guide the model's prediction.}

\cleardoublepage

\section{Prior work / Literature review}

\noindent

This thesis places itself at the intersection of computer vision, digital humanities, and epigraphy (the study of inscriptions).
Since it is a thesis of engineering, computer vision will remain the primary focus, both in terms of the research questions, and the methods used to answer them.
As a result, the other two fields will merely provide the context and motivation for the research.
The prior work and initial literature review in focus will thus concentrate on: 1. the relevant theory and state-of-the-art within computer vision (\acs{SAM}, etc.),
2. the palæographical and digital humanities work that this project will seek to support (SigLA),
and 3. the overall context thereof (Linear A).

\subsection{Research context}

For the context of Linear A and its corpus, the background section above can be referred to, citing the sources for: its size \parencite{Salgarella2020} and \parencite{Salgarella2025}, its new addition \parencite{RILA-S1}, and where it's accessed \parencite{EtudesCretoises21}.
And for the epigraphical and digital humanities work, the background section can also be referred to, citing the sources for SigLA: its paper \parencite{SalgarellaCastellanSigLA2020}, its about page \parencite{SigLAAboutPage}, and its document repository \parencite{SigLASearchDocuments}.
Here, additional context can also be provided by Ester Salgarella herself, as she is a co-supervisor of this project, and the main person behind SigLA.

\subsection{Computer vision literature}

As for the relevant theory and state-of-the-art within computer vision, the literature there can be further split into three parts;
1. sources of general image processing,
2. sources concerning SAM and its implementation,
and 3. related computer vision script segmentation work for comparison.

For image processing, a general theory book \parencite{DigitalImageProcessing} can be cited, as well as more method-specific papers (when other methods are used) like Otsu thresholding \parencite{Otsu1979}, and of course the OpenCV Python documentation \parencite{opencv_imgproc_docs} for short implementation reference.
Then, for \acs{SAM}, these are: its paper \parencite{sam}, its lighter implementation MobileSAM \parencite{mobile_sam} - and the newest version of that one MobileSAMv2 \parencite{mobile_sam_v2}.
Lastly, for related segmentation work there exists a paper using deep learning (in contrast to the methodology of this thesis) obtaining results better than state-of-the-art \parencite{DeepLearningSegmentation}, and a lighter EU-funded historical document segmentation tool \textit{dhSegment} \parencite{dhsegment} also using deep learning.


\section{Problem statement}
As previously stated, the aim of this project is to investigate whether parts of the process of creating annotated drawings for SigLA can be automated.
More specifically, whether some kind of semi-automatic segmentation tool using interactive and generalized computer vision - specifically the \acf{SAM} - can accelerate the workflow, while maintaining high-quality outputs that are still compatible with SigLA’s database.
But what would be the relative segmentation quality of such a tool? How automatic would it be exactly? And by how much could it improve the efficiency of the workflow?
For a decomposition of the problem statement into its three research questions - each accompanied by their respective operationalization - see the following subsection:

\subsection{Research questions}
\begin{enumerate}
\item \textbf{Segmentation performance}: 
How well does a \acs{SAM}-based segmentation approach perform quantitatively in terms of segmentation accuracy, compared to simple classical image processing baselines when applied to Linear A document images?
\note{Operationalization: \acs{SAM}-based masks will be compared to classical image-processing baselines using overlap metrics (e.g. \acs{IoU}/\acs{Dice}) against the ground truth segmentations derived from SigLA's Krita files for a minimum of 25 documents.}
\item \textbf{Degree of automation}: 
To which extent can user interaction be reduced by a \acs{SAM}-based semi-automatic segmentation workflow, while still maintaining high-quality outputs?
\note{Operationalization: The number of user interactions (e.g. clicks) required to achieve a certain level of segmentation quality (e.g. \acs{IoU}/\acs{Dice}) will be measured against ground truth (SigLA's Krita files) for a minimum of 25 documents.}
\item \textbf{Workflow efficiency}: 
How much reduction in annotation time can be achieved using the proposed tool compared to fully manual digitization, under controlled experimental conditions?
\note{Operationalization: The time it takes Ester Salgarella to annotate manually compared to being assisted by the proposed segmentation tool will be measured for a minimum of 5 documents.}
\end{enumerate}

\section{Empirical considerations}
For the evaluation of the solution, a dataset will be compiled together with Ester Salgarella of at least 25 images of original 
documents from the corpus, selected and cropped from \acs{Cefael}.
These images are subject to \acf{IPR} issues, and can thus not be displayed publicly if not through \acs{Cefael}.
However, they can still be used for the development and evaluation of the tool (and can be bought if necessary for show-casing within the thesis).
For direct reference to the primary source of data, the specific links to the \acs{Cefael} online webpages for each document used might have to be cited.

As ground truth, their corresponding digitized layered drawings, stored as Krita (.kra) files on SigLA (provided by Ester Salgarella) will be used and aligned for comparison with tool outputs.
The ground truth data might then have to be cleared of any marked 
“deleted signs” or other annotation elements that our model is not intended to reproduce.

The solution will then be built as a pipeline, consisting of
pre-processing, semi-automatic segmentation, post-processing, and export (e.g. as Krita files).
At last, the primary deliverable of this project will be delivered to the external partner and co-supervisor Ester Salgarella in the form of a proof-of-concept semi-automatic annotation tool.
All supporting components, 
including evaluation results, and analysis will be documented and presented in the written 
thesis. 

\section{Impact -- innovation and application}

The work of this thesis will present itself as a first proof-of-concept for the use of modern computer vision to accelerate the expansion of the SigLA database, supporting the further digitization of the Linear A corpus.
Its main contribution will be a semi-automatic annotation tool, which might prove useful to epigraphers for generating segmentations of original document images, as they can then be completed and refined more easily.
This approach could potentially also be applied to Linear B, as it is the whole field of Aegean studies that seem to suffer from severe under-digitization of the primary evidence, which is a huge problem for the scientific investigations of scholars.
Hopefully, it will inspire further research and development within the field of digital epigraphy, particularly in the application of generalized and interactive computer vision approaches such as \acs{SAM} on Linear A, and perhaps even in the broader context of use on small corpus ancient writing systems where the lack of big data is ill fit for approaches using deep learning.
